\section{A short imaginary-time slab}\label{sec:slab}

The next estimates describe the part of the Gibbs path measure that
will be changed later. Set $\delta=s^2$. For sufficiently small
density, $\beta>2\delta$, because
$\delta\leq4K^2\rho^{-1}$ and $\beta\geq\rho^{-2}$.
Let $P_{\mathrm{sym}}$ be the orthogonal symmetrization projector
on labeled coordinates. Write $F(Y,S)$ for the kernel, from $Y$
to $S$, of $e^{-(\beta-\delta)H}P_{\mathrm{sym}}$.
It is nonnegative and invariant under separate permutations of
each tuple. The inner slab has labeled paths $\omega_i$ from
$S_i$ to $E_i$, with unnormalized free torus bridge measures
$\mathsf K_{S_i,E_i}$. Put
\[
 \mathcal V_{ij}(\omega)=\int_0^\delta
 v(\omega_i(t)-\omega_j(t))\,dt,\qquad
 W(S,E;d\omega)=e^{-\sum_{i<j}\mathcal V_{ij}(\omega)}
 \prod_{i=1}^N\mathsf K_{S_i,E_i}(d\omega_i).
\]
In the closed measure $E=Y$, and the joint probability measure is
\begin{equation}
 Z_N^{-1}F(Y,S)\,dY\,dS\,W(S,Y;d\omega).
 \tag{14}\label{eq:14}
\end{equation}
Denote this closed probability law by $P_{\mathrm{cl}}$ and its
expectation by $\mathbb E_{\mathrm{cl}}$.
The normalization is the semigroup trace identity; the canonical
permutation and one-particle conventions agree with
\cite[Appendix~A, (A.1), (A.8)--(A.10), and (A.12)--(A.18)]{Ueltschi2006FeynmanCycles}. At each fixed
time, permutation-invariant position tests have the equilibrium
law $P$: insert the test as a multiplication operator, split the
semigroup at that time, and use trace cyclicity.

For clarity, a torus bridge includes its lift choice. Fix a
representative of its end, lift its start in every possible way,
and assign a displacement $d_\ell$ to that lift. Its weight is
$k(d_\ell)=(4\pi\delta)^{-3/2}e^{-|d_\ell|^2/(4\delta)}$
times the law of a Euclidean bridge with generator $\Delta$,
projected to the torus. This gives the unnormalized torus bridge
after summation over $\ell\in\mathbb Z^3$. All lifted tests below
refer to this lift, not to a pointwise choice of nearest torus
representatives along the trajectory.

Fix a large number $D$, put $R=D+30$ and $R_A=2R+20$, and later
choose $0<\varepsilon<1/2$. The short-distance scale is
$r=\zeta\rho^{-1/3}$ from Lemma~\ref{lem:isolation}. The order of
choices allows $\zeta$ to depend on $D,\varepsilon$, followed by
a small-density restriction. In particular $r\to\infty$,
$r/s\to0$, and $2r<\delta$.
For an anchor $a$, define
\begin{equation}
 \begin{aligned}
 \mathcal A^-(a;\omega)&=\int_0^{\delta/2}w(t)
 \sum_i e^{-d(\omega_i(t),a)^2/[1000(1+t)]}
 \mathbf1_{\{d(\omega_i(t),a)\leq R_As\}}\,dt,\\
 &\hspace{15mm}w(t)=\min(1,t^{-3/2}),\qquad w(0)=1.
 \end{aligned}
 \tag{15}\label{eq:15}
\end{equation}
Define $\mathcal A^+$ by replacing $\omega_i(t)$ by
$\omega_i(\delta-t)$. A subcollection in place of $\omega$ means
that only its paths occur in the sum.

A closed slot $i$ is \emph{eligible} if all of the following hold:
its start $S_i$ and end $Y_i$ are each at distance greater than $r$
from the other points of the respective tuple; its lifted path
stays within $Ds$ of its end; in the first and last time intervals
of length $r$ it stays within $r/10$ of its respective start and
end; and both $\mathcal A^-(S_i;\omega)$ and
$\mathcal A^+(Y_i;\omega)$ are at most a fixed threshold $C_0$.
A cell is \emph{usable} if its end count belongs to
$[b_-m,b_+m]$ and at least a proportion $1-\varepsilon$ of those
slots are eligible.
For an individual bridge, $\mathrm{cut}_R$ denotes just the lifted
macro condition with radius $Rs$ and the same two interval
conditions with radius $r/10$. In this definition every condition
uses that bridge's own endpoints.

\begin{lemma}[Single-bridge estimates]\label{lem:single-bridge}
Suppose a bridge's endpoints have torus distance at most
$(D+10)s$. When $L$ is sufficiently large compared with $Rs$,
$\mathrm{cut}_R$ permits only the nearest lift. Under the probability
bridge $\pi_{S,E}$ for that lift, its probability is at least a
constant $p_*>0$. Under the same unconditioned probability bridge
$\pi_{S,E}$, the expected interaction with any fixed trajectory
$\omega_j$ is at most $C_v$, and the contribution
of times $[r,\delta-r]$ is at most $C_v/\sqrt r$.
For any fixed collection $\omega_{\mathrm{fix}}$,
\begin{equation}
 \mathbb E_{\pi_{S,E}}\left[
 \mathbf1_{\mathrm{cut}_R}
 \sum_{j\in\mathrm{fix}}\int_0^\delta
 v(\omega(t)-\omega_j(t))\,dt\right]
 \leq C(v,D)\bigl[
 \mathcal A^-(S;\omega_{\mathrm{fix}})+
 \mathcal A^+(E;\omega_{\mathrm{fix}})\bigr].
 \tag{16}\label{eq:16}
\end{equation}
Finally, for any fixed collection of at most $N$ paths and a
uniform torus point $x$,
$\mathbb E_x\mathcal A^\pm(x;\omega)\leq CK^2$.
These bounds are independent of temperature, density, and volume
after the stated choices; $p_*$ may be fixed after $D$.
\end{lemma}

\begin{proof}
Take $L>4Rs$. Any nonnearest lift has displacement greater than
$Rs$, so it fails the cut already at time zero. In the nearest
lift write the bridge as its straight interpolation plus
$B_t-(t/\delta)B_\delta$, where $B$ has generator $\Delta$.
By Brownian scaling the event that this fluctuation has supremum
at most $s$ has a fixed positive probability; on it the macro
distance is at most $(D+11)s<Rs$.
The deterministic drift during the first interval is at most
$(D+10)r/s$, which is at most $r/20$ once $s\geq20(D+10)$.
The Brownian maximal bound and the Gaussian tail of $B_\delta$
show that the remaining fluctuation exceeds $r/20$ on this
interval with probability at most $Ce^{-cr}$: separately bound
$\sup_{t\leq r}|B_t|$ by $r/40$ and
$|B_\delta|$ by $\delta/40$. The same argument applies at the
other end by reversal. Subtracting these two small failure
probabilities from the macro event proves the lower bound $p_*$.

At time $t$ the nearest-lift bridge is Gaussian with coordinate
variance $2t(\delta-t)/\delta$. For
$t'=\min(t,\delta-t)$ this variance lies in $[t',2t']$.
Its periodic density is bounded by $C(t')^{-3/2}$ when
$t'\leq\delta\ll L^2$. Thus the expectation of $v$ against any
fixed point is at most $C_vw(t')$: use $\|v\|_\infty$ for
$t'\leq1$ and the density bound with $\|v\|_1$ for $t'\geq1$.
Integration gives a bound $C_v$ on the full interaction and
$C_v\int_r^\infty t^{-3/2}\,dt\leq C_v/\sqrt r$ on the
middle interval.

To retain spatial decay, first take $t\leq\delta/2$ and use
the start $S$ as anchor. The interpolation center is within
$(D+10)t/s\leq(D+10)\sqrt t$ of $S$.
For $t\geq1$, the Gaussian density integrated over the periodically
lifted range balls around a fixed point $y$ is at most
\[
 C(v,D)t^{-3/2}\sum_{\ell\in\mathbb Z^3}
 e^{-|y+\ell L-S|^2/(16t)}.
\]
Indeed $|u+b|^2\geq|u|^2/2-|b|^2$ absorbs the center displacement
and the range radius, and $|b|\leq R_v+(D+10)\sqrt t$.
Extracting $e^{-d(y,S)^2/(32t)}$ leaves a uniformly bounded
periodic Gaussian sum, because $t\leq\delta\ll L^2$.
For $0<t\leq1$, use probability at most one within a fixed
distance $2(R_v+D+10)$ of $S$. Beyond this distance, a hit
requires a Gaussian fluctuation at least $d(y,S)/2$; its tail
has the same, or stronger, spatial decay after increasing
$C(v,D)$. We have thus bounded the expected interaction density
by
$C(v,D)w(t)e^{-d(y,S)^2/[1000(1+t)]}$.

On the cut, a possible hit also requires
$d(y,S)\leq (R+D+10)s+R_v<R_As$ for sufficiently small density.
This follows by bounding the bridge's distance from the end by
$Rs$ and the endpoint displacement by $(D+10)s$.
Consequently the indicator in (15) may be inserted in this upper
bound. Reverse time for the other half, sum over the fixed
trajectories, and integrate to obtain (16).
Finally the integral over a uniform anchor of the Gaussian factor
is at most $C(1+t)^{3/2}/V$. Since
$w(t)(1+t)^{3/2}\leq C$, the mean score is at most
$CN\delta/V\leq C\rho s^2\leq CK^2$.
\end{proof}
